\documentclass[11pt,a4paper]{article}

% ============================================================
% FORM-TL-05 -- Informe de Revision Tecnica Cruzada
% Equipo revisor: AGLS   |   Equipo/PFC revisado: BCEL - SGA Escuela Provincias Unidas
% ============================================================

\usepackage[spanish,es-tabla]{babel}
\usepackage[utf8]{inputenc}
\usepackage[T1]{fontenc}
\usepackage[margin=2.5cm]{geometry}
\usepackage[expansion=false]{microtype}
\emergencystretch=3em
\sloppy
\usepackage{booktabs}
\usepackage{tabularx}
\usepackage{xltabular}
\usepackage[most]{tcolorbox}
\usepackage{listings}
\usepackage{xcolor}
\usepackage{xurl}
\usepackage{enumitem}
\usepackage{tikz}
\usetikzlibrary{arrows.meta,positioning,fit,backgrounds}
\usepackage{fancyhdr}
\usepackage{titlesec}
\usepackage{longtable}
\usepackage[backend=biber,style=ieee,sorting=none]{biblatex}
\addbibresource{referencias.bib}
\usepackage[hidelinks,breaklinks]{hyperref}

\definecolor{uteqblue}{RGB}{36,58,118}
\definecolor{boxgray}{RGB}{245,246,248}
\definecolor{critred}{RGB}{178,34,34}
\definecolor{majorange}{RGB}{204,120,0}
\definecolor{minoryellow}{RGB}{150,120,0}
\definecolor{sugblue}{RGB}{70,110,160}

\titleformat{\section}{\large\bfseries\color{uteqblue}}{\thesection}{0.6em}{}
\titleformat{\subsection}{\normalsize\bfseries\color{uteqblue}}{\thesubsection}{0.6em}{}

\lstset{
    basicstyle=\ttfamily\footnotesize,
    breaklines=true,
    frame=single,
    numbers=left,
    numberstyle=\tiny\color{gray},
    xleftmargin=1.8em,
    backgroundcolor=\color{boxgray},
    columns=fullflexible,
    showstringspaces=false
}

\newtcolorbox{sevcrit}{colback=critred!8,colframe=critred,fonttitle=\bfseries,title=Severidad: Cr\'itica}
\newtcolorbox{sevmajor}{colback=majorange!8,colframe=majorange,fonttitle=\bfseries,title=Severidad: Mayor}
\newtcolorbox{sevminor}{colback=minoryellow!8,colframe=minoryellow,fonttitle=\bfseries,title=Severidad: Menor}
\newtcolorbox{sevsug}{colback=sugblue!8,colframe=sugblue,fonttitle=\bfseries,title=Sugerencia}
\newtcolorbox{notebox}[1][]{colback=uteqblue!5,colframe=uteqblue,fonttitle=\bfseries,#1}

\pagestyle{fancy}
\fancyhf{}
\lhead{\small UTEQ \textperiodcentered\ ISR-701 \textperiodcentered\ FORM-TL-05}
\rhead{\small Revisi\'on cruzada AGLS $\to$ BCEL}
\cfoot{\small\thepage}

\title{
    \vspace{-1.5cm}
    {\small UNIVERSIDAD T\'ECNICA ESTATAL DE QUEVEDO\\Facultad de Ciencias de la Computaci\'on\\Carrera de Ingenier\'ia de Software}\\[0.6cm]
    \textbf{\Large Informe de Revisi\'on T\'ecnica Cruzada del C\'odigo del PFC}\\[0.2cm]
    \large FORM-TL-05 --- Principios SOLID, patrones de dise\~no, separaci\'on de capas,\\
    manejo de excepciones distribuidas y observabilidad\\[0.3cm]
    \normalsize Sistema auditado: \textbf{SGA -- Sistema de Gesti\'on Acad\'emica Distribuido}\\
    (equipo autor \textbf{BCEL}, sigla institucional ``Escuela Provincias Unidas'')
}
\author{Equipo revisor: \textbf{AGLS}\\ \small (PFC propio del equipo revisor: TiendaTech -- comercio electr\'onico, no auditado en este informe)}
\date{Per\'iodo 2026--2027 \textperiodcentered\ Semana 16}

\begin{document}

    \maketitle
    \thispagestyle{empty}

    \begin{center}
        \renewcommand{\arraystretch}{1.3}
        \begin{tabularx}{\textwidth}{@{}>{\raggedright\arraybackslash}p{4.2cm}X@{}}
\toprule
\multicolumn{2}{l}{\textbf{Datos del equipo REVISOR (AGLS)}}\\
\midrule
Integrante 1 -- Rol & Jhinson Stalyn Aucatoma Celorio -- Bloques A, B y C (SOLID, patrones y separaci\'on de capas) \\
Integrante 2 -- Rol & Andy Paul S\'anchez Pilaloa -- Bloques D y E (excepciones distribuidas y trazabilidad) \\
\midrule
\multicolumn{2}{l}{\textbf{Datos del PFC REVISADO (BCEL)}}\\
\midrule
Equipo revisado (sigla) & BCEL \\
Nombre del sistema & SGA -- Sistema de Gesti\'on Acad\'emica Distribuido (Escuela Provincias Unidas) \\
URL del repositorio revisado & \small\url{https://github.com/LEO23as/sga-sistema-distribuido} \\
Commit auditado (SHA corto) & \texttt{\small ff6f8b0} \small (completo: \texttt{ff6f8b0bf3f4200ae4deb9d6d7833e8bce7330f4}) \\
Fecha y hora de la auditor\'ia & 11 de agosto de 2026, 08:30 \\
\bottomrule
        \end{tabularx}
    \end{center}

    \vspace{0.4cm}
    \begin{notebox}[title=Nota de verificaci\'on del objeto revisado \label{sec:verificacion}]
        Siguiendo la pr\'actica profesional de revisi\'on de c\'odigo moderna, cuyo valor documentado excede la sola detecci\'on de defectos y se centra en la transferencia de conocimiento entre equipos~\cite{bacchelli2013,sadowski2018}, el an\'alisis inicial se elabor\'o sobre el archivo \texttt{sga-sistema-distribuido-main.zip} entregado como material de trabajo, correspondiente a una descarga ``Download ZIP'' de la rama \texttt{main} y sin directorio \texttt{.git}. Posteriormente el equipo AGLS complet\'o la verificaci\'on exigida en la Secci\'on 6 de la gu\'ia FORM-TL-05 sobre un clon real del repositorio:
        \begin{lstlisting}[language=bash,numbers=none]
git clone https://github.com/LEO23as/sga-sistema-distribuido.git revision-pfc
cd revision-pfc
git log -1 --format="%H | %aI | %an"
git rev-parse --short HEAD
git ls-files | wc -l
        \end{lstlisting}
        confirmando como commit auditado el \texttt{ff6f8b0} (SHA completo \texttt{ff6f8b0bf3f4200ae4deb9d6d7833e8bce7330f4}), verificado el 11 de agosto de 2026 a las 08:30, valores ya reflejados en la portada. Todos los hallazgos de este informe referencian rutas de archivo estables dentro del \'arbol de trabajo (\texttt{sga-principal/}, \texttt{microservicio-docente/}, \texttt{microservicio-secretaria/}, \texttt{microservicio-soporte/}), por lo que siguen siendo v\'alidos independientemente de cambios posteriores al commit auditado, siempre que se referencien contra ese SHA exacto.
    \end{notebox}

    \newpage
    \tableofcontents
    \thispagestyle{empty}
    \newpage
    \setcounter{page}{1}

================================================================
\section{Mapa de arquitectura reconstruido}
\label{sec:arquitectura}

El sistema BCEL se compone de un servicio n\'ucleo (\texttt{sga-principal}, Java 17 / Spring Boot) y tres microservicios aut\'onomos (\texttt{microservicio-docente} en Django REST + gRPC, y \texttt{microservicio-secretaria} / \texttt{microservicio-soporte}, ambos Spring Boot pese a que el \texttt{README.md} los documenta como ``Node.js/Express''\footnote{\texttt{README.md}, tabla de arquitectura: declara Node.js/Express para Secretar\'ia y Soporte; el c\'odigo real en \texttt{microservicio-secretaria/backend/pom.xml} y \texttt{microservicio-soporte/backend/pom.xml} es Maven/Spring Boot. Documentaci\'on desactualizada, ver Anexo de sugerencias.}). La reconstrucci\'on se hizo a partir de los \texttt{application.properties}/\texttt{settings.py}, los clientes gRPC y \texttt{docker-compose.yml}, no de la documentaci\'on.

\begin{center}
\resizebox{\textwidth}{!}{%
\begin{tikzpicture}[
  node distance=0.9cm and 1.3cm,
  box/.style={draw, rounded corners, minimum width=2.9cm, minimum height=1cm, align=center, font=\footnotesize},
  svc/.style={box, fill=uteqblue!10},
  db/.style={box, fill=orange!15, minimum width=4.6cm},
  gw/.style={box, fill=gray!15},
  fe/.style={box, fill=green!10},
  arr/.style={-{Latex[length=2mm]}, thick},
  grpcarr/.style={-{Latex[length=2mm]}, thick, dashed},
]
\node[fe] (front) {Frontends React\\(5173/5174/5175/5176)};
\node[gw, below=of front] (haproxy) {HAProxy\\(solo balancea\\sga-principal: 8080/9092)};
\node[svc, below=of haproxy] (principal) {sga-principal\\Spring Boot\\REST 8080 / gRPC 9092};
\node[svc, below left=of principal, xshift=-0.9cm] (docente) {microservicio-docente\\Django REST + gRPC\\8081 / 9091};
\node[svc, below right=of principal, xshift=0.9cm] (secretaria) {microservicio-secretaria\\Spring Boot\\8082 / gRPC cliente};
\node[svc, right=1.6cm of secretaria] (soporte) {microservicio-soporte\\Spring Boot + etcd\\(elecci\'on de l\'ider)\\8083 / gRPC cliente};
\node[db, below=1.6cm of principal] (pg) {PostgreSQL \'unico -- 3.23.195.43:5433\\esquemas: sga\_principal, sga\_docente, public\\(search\_path cruzado, ver H05)};

\draw[arr] (front) -- (haproxy);
\draw[arr] (haproxy) -- (principal);
\draw[grpcarr] (principal) -- node[left,font=\scriptsize]{gRPC :9091} (docente);
\draw[grpcarr] (secretaria) -- node[right,font=\scriptsize]{gRPC :9092} (principal);
\draw[grpcarr] (soporte) -- node[below,font=\scriptsize]{gRPC :9092} (principal);
\draw[arr,dotted] (front) -- (secretaria);
\draw[arr,dotted] (front) -- (soporte);
\draw[arr,dotted] (front) -- (docente);
\draw[arr] (principal) -- (pg);
\draw[arr] (docente) -- (pg);
\draw[arr] (secretaria) -- (pg);
\draw[arr] (soporte) -- (pg);
\end{tikzpicture}%
}
\end{center}
\footnotesize
Flechas s\'olidas: REST/HTTP. Flechas discontinuas: gRPC (\texttt{PLAINTEXT}, sin TLS). Flechas punteadas: acceso directo de cada frontend a su propio backend, sin pasar por HAProxy (HAProxy s\'olo tiene \texttt{frontend}/\texttt{backend} definidos para \texttt{sga-principal}, ver \texttt{infra/haproxy/haproxy.cfg} l\'ineas 24--47). \texttt{sga-principal} act\'ua de facto como \emph{hub} gRPC: \texttt{secretaria} y \texttt{soporte} lo consultan para el cat\'alogo institucional (a\~nos, grados, estudiantes), y \'el a su vez consulta a \texttt{docente} para actividades/asistencias. \texttt{soporte} incorpora adem\'as un algoritmo propio de elecci\'on de l\'ider sobre \texttt{etcd} (\texttt{election/LeaderElectionService.java}), un acierto de dise\~no distribuido que se reconoce en la Secci\'on~\ref{sec:conclusiones}.
\normalsize

%
    ================================================================
    \section{Tabla consolidada de hallazgos}
    \label{sec:hallazgos}

    Cada hallazgo se clasifica adem\'as por el atributo de calidad de producto que compromete, siguiendo la nomenclatura de ISO/IEC~25010~\cite{iso25010} (mantenibilidad, fiabilidad, seguridad o eficiencia de desempe\~no, seg\'un el caso); esa clasificaci\'on se detalla en el cuerpo de cada issue del Anexo A (archivo \texttt{issues\_plan.md}) para no duplicar columnas en esta tabla. Cada hallazgo referencia archivo y l\'inea dentro del \'arbol de trabajo del repositorio revisado. La columna \emph{Item} usa la nomenclatura de la Secci\'on 5 de la gu\'ia FORM-TL-05 (bloques A--E). La columna \emph{Issue} enlaza al n\'umero que deber\'a asignarse al abrir el issue en el repositorio de BCEL (Secci\'on~\ref{sec:issues}).

        {\small
    \begin{xltabular}{\textwidth}{@{}p{0.55cm}p{0.95cm}X p{3.5cm}p{1.05cm}p{0.6cm}@{}}
        \toprule
        \textbf{ID} & \textbf{Item} & \textbf{Hallazgo} & \textbf{Evidencia (ruta:l\'inea)} & \textbf{Sever.} & \textbf{\#} \\
        \midrule
        \endhead
        H01 & A1 & \texttt{ImportacionExcelService} concentra parseo de CSV, Excel y PDF y la persistencia de estudiantes/matr\'iculas en una sola clase de 388 l\'ineas. & {\raggedright\scriptsize\url{sga-principal/.../service/ImportacionExcelService.java:43-388}\par} & Mayor & 1 \\
        H02 & A2 & Agregar un formato de archivo nuevo (p.\,ej. \texttt{.ods}) obliga a editar el \texttt{if/else} existente en \texttt{parsearArchivo} y a\~nadir un m\'etodo \texttt{leerX} paralelo. & {\raggedright\scriptsize\url{.../ImportacionExcelService.java:61-70}\par} & Menor & 2 \\
        H03 & A5/C2 & \texttt{AuthController} inyecta \texttt{UsuarioRepository} directamente y lo consulta con \texttt{.orElseThrow()} sin manejar, saltando la capa de servicio. & {\raggedright\scriptsize\url{.../controller/AuthController.java:9,23,39-41}\par} & Mayor & 3 \\
        H04 & A1/A5/C2 & \texttt{PrincipalGrpcService} (adaptador gRPC, capa de comunicaci\'on) inyecta 4 repositorios adem\'as de 2 servicios y mezcla marshalling, acceso a datos y reglas de 5 agregados distintos en 382 l\'ineas. & {\raggedright\scriptsize\url{.../grpc/PrincipalGrpcService.java:37-44}\par} & Mayor & 4 \\
        H05 & B5/C5 & El esquema de \texttt{microservicio-docente} declara \texttt{search\_path=sga\_docente,sga\_principal,public}: puede leer/escribir directamente el esquema de otro servicio. Los 4 backends adem\'as convergen en la misma instancia f\'isica de PostgreSQL. & {\raggedright\scriptsize\url{micro_docente/settings.py:53-65} (l.63)\par} & Cr\'itica & 5 \\
        H06 & C5 & Contrase\~na de BD, secreto JWT, contrase\~na de aplicaci\'on de Gmail y token interno de gRPC hardcodeados y versionados en texto claro; credenciales reales de acceso publicadas en el \texttt{README.md}. & {\raggedright\scriptsize\url{application.properties:10,29,33-34};\newline \url{settings.py:58};\newline \url{InternalAuthInterceptor.java:15}; README.md\par} & Cr\'itica & 6 \\
        H07 & D2 & Ning\'un cliente gRPC del sistema (\texttt{Actividad}/\texttt{Asistencia}/\texttt{DocenteGrpcClient}, \texttt{PrincipalGrpcClient}, \texttt{TecnicoGrpcClient}) fija \texttt{withDeadlineAfter}; las llamadas son bloqueantes sin l\'imite de espera propio. & {\raggedright\scriptsize\url{.../grpc/ActividadGrpcClient.java:30-37} (5 clientes)\par} & Mayor & 7 \\
        H08 & D3/D4 & No existe pol\'itica de reintentos ni respuesta de degradaci\'on: los \texttt{catch} s\'olo relanzan la excepci\'on original o la traducen 1:1 a un error HTTP, sin valor de reserva. & {\raggedright\scriptsize\url{ActividadGrpcClient.java:33-36,42-45,52-55,67-70,77-80}\par} & Mayor & 7 \\
        H09 & D5 & Formato de error no uniforme entre servicios: \texttt{secretaria} y \texttt{soporte} tienen \texttt{@RestControllerAdvice} propio; \texttt{sga-principal} no tiene ning\'un manejador global y depende de \texttt{ResponseStatusException} ad-hoc por m\'etodo. & {\raggedright\scriptsize \texttt{sga-principal}: sin \texttt{@ControllerAdvice}\newline (0 result.) vs.\newline \url{.../soporte/.../GlobalExceptionHandler.java:1-127}\par} & Mayor & 8 \\
        H10 & E1 & Registro por consola (\texttt{System.out}/\texttt{err.println}) en vez de logger inyectado, incluyendo el punto de entrada de la aplicaci\'on. & {\raggedright\scriptsize\url{DocenteContextController.java:31-86};\newline \url{SgaPrincipalApplication.java:14};\newline \url{DataSeedRunner.java:47-182};\newline \url{ActividadGrpcClient.java:34,43,53,68,78}\par} & Menor & 9 \\
        H11 & E2 & No existe identificador de correlaci\'on en ning\'un punto del sistema (0 coincidencias de \texttt{MDC}/\texttt{correlationId}/\texttt{traceId}/\texttt{X-Request-Id} en los 3 backends Java ni en \texttt{settings.py} de Django). & {\raggedright\scriptsize B\'usqueda global \texttt{grep -rln} en los 4 servicios (0 result.)\par} & Mayor & 10 \\
        H12 & E4 & \texttt{main()} imprime en cada arranque el hash bcrypt de una contrase\~na hardcodeada; el \texttt{README.md} publica usuario y contrase\~na reales de administrador y docente en texto claro. & {\raggedright\scriptsize\url{SgaPrincipalApplication.java:12-16}; README.md (\S\ Credenciales)\par} & Cr\'itica & 11 \\
        H13 & B4 & La \'unica puerta de enlace (HAProxy) balancea exclusivamente \texttt{sga-principal}; \texttt{docente}, \texttt{secretaria} y \texttt{soporte} no tienen \texttt{backend} definido, por lo que el frontend debe conocer y llamar cada puerto de servicio directamente. & {\raggedright\scriptsize\url{infra/haproxy/haproxy.cfg:24-47}\par} & Menor & 12 \\
        H14 & Sug. & Higiene de repositorio: clase p\'ublica vac\'ia con nombre en min\'uscula (\texttt{repository.java}); binarios \texttt{protoc-*.exe} versionados en \texttt{target/} pese a que \texttt{.gitignore} excluye \texttt{**/target/}; \texttt{.gitignore} con codificaci\'on corrupta en sus \'ultimas l\'ineas. & {\raggedright\scriptsize\url{.../repository/repository.java};\newline \url{.../soporte/backend/target/protoc-plugins/}; .gitignore\par} & Sug. & 13 \\
        \bottomrule
    \end{xltabular}
    }

\end{document}
